\begin{helper}[Bennett entropy lower bound]
For every real $z\ge0$,
\[
 (1+z)\log(1+z)-z\ge\frac{z^2}{2(1+z/3)}.
\]
\end{helper}
\begin{proof}
On $[0,\infty)$ the quantities $1+x$, $x+3$, and $1+x/3$ are
strictly positive. The identity
\[
 \frac{x^2}{2(1+x/3)}=\frac32\left(x-3+\frac9{x+3}\right)
\]
follows by multiplying by the positive denominators. Define the
difference between the left entropy and this rational expression by
\[
 f(x)=(1+x)\log(1+x)-x-\frac32\left(x-3+\frac9{x+3}\right).
\]
It is differentiable at every $x\ge0$, with
\[
 f'(x)=g(x)=\log(1+x)-\frac32+\frac{27}{2(x+3)^2},\qquad
 g'(x)=\frac1{1+x}-\frac{27}{(x+3)^3}.
\]
Both denominators in $g'$ are positive. Consequently $g'(x)\ge0$
is equivalent to $(x+3)^3\ge27(1+x)$. Their difference is
$x^3+9x^2\ge0$ for $x\ge0$, which proves this sign. The functions
$f$ and $g$ are continuous on $[0,\infty)$, so the mean value theorem
on each compact subinterval shows first that $g$ is nondecreasing.
Since $g(0)=0$, this gives $f'(x)=g(x)\ge0$ on the same half-line.
Applying the mean value theorem once more makes $f$ nondecreasing.
Finally $f(0)=0$ implies $f(z)\ge0$ for every $z\ge0$. Substituting
the displayed rational identity gives the result.
\end{proof}
